% UTF-8 / pdfLaTeX source.
% Compile example:
%   pdflatex fmo_matrix_projection_explained.tex

\documentclass[11pt]{article}

\usepackage[T1]{fontenc}
\usepackage[utf8]{inputenc}
\usepackage{amsmath,amssymb}
\usepackage{geometry}
\usepackage{array}
\usepackage{longtable}

\geometry{a4paper,margin=25mm}

\newcommand{\Hcore}{\mathbf H}
\newcommand{\Venv}{\mathbf V}
\newcommand{\Smat}{\mathbf S}
\newcommand{\Amat}{\mathbf A}
\newcommand{\Bmat}{\mathbf B}
\newcommand{\Tmat}{\mathbf T}
\newcommand{\Cmat}{\mathbf C}
\newcommand{\Cr}{\mathbf C_{\mathrm r}}
\newcommand{\Fmat}{\mathbf F}
\newcommand{\Fr}{\mathbf F_{\mathrm r}}
\newcommand{\Sr}{\mathbf S_{\mathrm r}}
\newcommand{\Dmat}{\mathbf D}
\newcommand{\Dr}{\mathbf D_{\mathrm r}}
\newcommand{\Gmat}{\mathbf G}

\title{FMO Boundary Treatment and the Removal of Hybrid Orbitals\\
by a Rectangular Transformation}
\author{}
\date{Created: July 5, 2026}

\begin{document}
\maketitle

\section*{Purpose of this document}

When a molecule is divided into fragments in the FMO method, a
``forbidden hybrid orbital'' pointing along a severed bond remains at the
cut. This document describes how to remove that orbital from the
calculation.

Stating the conclusion first: there is no need to push the orbital up
with a large level shift \(B\), as is done conventionally. Instead, one
can solve the SCF from the outset in a smaller space that does not
contain that direction. We call this operation ``deflation by a
rectangular transformation.''

\section{What is the problem in the first place?}

FMO stands for Fragment Molecular Orbital. It is a method that performs
quantum chemical calculations by dividing a large molecule into several
small fragments. For systems such as proteins and polymers, which are
too heavy to treat as a whole, this idea becomes essential.

However, when a molecule is divided into fragments, covalent bonds may
be severed. For example, one may cut a carbon--carbon single bond in the
middle, assigning the left side to fragment \(x\) and the right side to
another fragment.

When a covalent bond is severed, the ``direction that originally formed
the bond'' remains at the cut point. If the carbon is \(sp^3\)
hybridized, the problem is the one hybrid orbital, out of the four
\(sp^3\) hybrids, that points along the severed bond.

Schematically, the situation is as follows.

\begin{verbatim}
Fragment x side                 Other fragment side

   allowed directions -- C_alpha ---- [ cut-direction sp3 hybrid ]
                                           ^
                                  not allowed in x calculation
\end{verbatim}

In this document we denote this ``cut-direction hybrid orbital'' by
\[
  \varphi_i^h .
\]
If there are several of them, we number them \(i=1,2,\ldots\)

If this orbital is left usable in the calculation of fragment \(x\),
the electrons of fragment \(x\) leak into a direction that properly
belongs to the partner fragment. This is physically undesirable.

Therefore, the FMO boundary treatment must do the following.

\begin{quote}
Remove the hybrid orbitals pointing along the severed bonds from the
SCF variational space of the fragment.
\end{quote}

Here, the ``variational space'' is the set of degrees of freedom that
the SCF is allowed to vary when searching for molecular orbitals. If an
unwanted direction remains there, the SCF may end up using it.

\section{A minimum of quantum chemistry vocabulary}

\subsection{AOs and MOs}

In quantum chemical calculations, molecular orbitals (MOs) are
expressed as linear combinations of atomic orbitals (AOs):
\[
  \psi_p = \sum_\mu C_{\mu p}\chi_\mu .
\]
Here \(\chi_\mu\) are the AOs, \(\psi_p\) are the MOs, and
\(C_{\mu p}\) are the coefficients specifying how much of each AO is
mixed into each MO.

The matrix collecting all MO coefficients is written \(\Cmat\).
Roughly speaking, the SCF is a procedure that adjusts \(\Cmat\) to
lower the electronic energy.

\subsection{AOs are, in general, not orthogonal}

The overlap between AOs is
\[
  S_{\mu\nu} = \langle \chi_\mu | \chi_\nu \rangle ,
\]
and the matrix collecting these is the overlap matrix \(\Smat\).

If the AOs were perfectly orthogonal we would have
\(\Smat=\mathbf I\). In practice, however, with atom-centered basis
sets, AOs on different atoms overlap in space. Hence, in general,
\[
  \Smat \ne \mathbf I .
\]
This fact---that AOs are non-orthogonal---is the reason \(\Smat\)
appears repeatedly in the equations below.

\subsection{The hybrid orbitals to be removed}

The cut-direction hybrid orbitals can also be expressed as linear
combinations of AOs:
\[
  \varphi_i^h = \sum_\mu A_{\mu i}\chi_\mu .
\]
We write the matrix whose columns are these coefficients as \(\Amat\).
If \(N\) is the number of AOs and \(m\) is the number of hybrid
orbitals to be removed, then \(\Amat\) is an \(N\times m\) matrix. For
example, if there is only one cut-direction \(sp^3\) hybrid, then
\(m=1\) and \(\Amat\) is a single column vector.

\section{What does it mean to ``not use the forbidden directions''?}

Whether an MO \(\psi_p\) contains the hybrid orbital \(\varphi_i^h\) to
be removed can be checked by their inner product:
\[
  \langle \varphi_i^h | \psi_p \rangle .
\]
If this vanishes, \(\psi_p\) has no component along \(\varphi_i^h\).

We want this to vanish for every removal target \(i\) and every MO
\(p\); in matrix form,
\[
  \Amat^\dagger \Smat \Cmat = \mathbf 0 .
\]

The reason \(\Smat\) enters here is that the AOs are non-orthogonal.
Indeed,
\begin{align*}
\langle \varphi_i^h | \psi_p \rangle
&=
\left\langle
\sum_\mu A_{\mu i}\chi_\mu
\middle|
\sum_\nu C_{\nu p}\chi_\nu
\right\rangle \\
&=
\sum_{\mu\nu} A_{\mu i}^{*}
\langle \chi_\mu | \chi_\nu \rangle
C_{\nu p} \\
&=
\sum_{\mu\nu} A_{\mu i}^{*} S_{\mu\nu} C_{\nu p}.
\end{align*}
Collecting this into matrix form gives \(\Amat^\dagger\Smat\Cmat\).

The mathematical goal of the boundary treatment can therefore be
summarized in a single line:
\[
  \boxed{\Amat^\dagger \Smat \Cmat = \mathbf 0}
\]
This states that ``no MO contains a cut-direction hybrid orbital.''

\section{The conventional method: pushing up with a level shift}

The conventional method assigns a very large energy to the forbidden
orbitals. The idea is simple:

\begin{quote}
Make using that direction so energetically costly that the SCF never
chooses it.
\end{quote}

This method is called level-shift projection, or hybrid orbital
projection.

With the core Hamiltonian \(\Hcore\) and the environmental
electrostatic potential \(\Venv\), the level-shifted matrix takes the
form
\[
  \widetilde{\Hcore}
  =
  \Hcore+\Venv+\Smat\Amat\Bmat\Amat^\dagger\Smat .
\]
Here \(\Bmat\) is a diagonal matrix of large positive values:
\[
  \Bmat = \mathrm{diag}(B_1,B_2,\ldots,B_m).
\]
Typically, very large values are used for \(B_i\). This raises the
orbital energies along the \(\varphi_i^h\) directions substantially,
with the expectation that they are expelled from the occupied space.

\subsection{Why the projection term has the form
\(\Smat\Amat\Bmat\Amat^\dagger\Smat\)}

The component form of the projection term is
\[
  P_{\mu\nu}
  =
  \sum_i B_i
  \langle \chi_\mu | \varphi_i^h \rangle
  \langle \varphi_i^h | \chi_\nu \rangle .
\]

The hybrid orbital \(\varphi_i^h\) is represented by the AO
coefficients \(\Amat\), but the AOs are non-orthogonal. Consequently,
\[
  \langle \chi_\mu | \varphi_i^h \rangle
  =
  (\Smat\Amat)_{\mu i} ,
\]
and an \(\Smat\) likewise emerges from the inner product on the right.
Therefore
\[
  \mathbf P = \Smat\Amat\Bmat\Amat^\dagger\Smat .
\]

If the AOs were orthogonal from the start, with \(\Smat=\mathbf I\),
the projection term would simply be
\[
  \Amat\Bmat\Amat^\dagger .
\]
In other words, the two factors of \(\Smat\) on the left and right are
a consequence of using non-orthogonal AOs.

\section{Weaknesses of the level-shift method}

Increasing \(B_i\) makes the components of the unwanted hybrid orbitals
smaller. With finite \(B_i\), however, they never vanish completely.
The residual components are roughly of size
\[
  O(B^{-1}) .
\]

On the other hand, making \(B_i\) too large places extremely large
values inside the matrix. The condition number then deteriorates and
the computation becomes numerically unstable. SCF convergence may also
suffer.

\begin{center}
\begin{tabular}{|p{0.22\linewidth}|p{0.34\linewidth}|p{0.34\linewidth}|}
\hline
\textbf{Choice of \(B\)} & \textbf{Advantage} & \textbf{Drawback} \\
\hline
Small & Numerically relatively stable & Unwanted components tend to remain \\
\hline
Large & Unwanted components become small & Condition number tends to deteriorate \\
\hline
Infinite & Ideally complete removal & Unusable in actual computation \\
\hline
\end{tabular}
\end{center}

To avoid this problem, a method that does not use \(B\) at all is
needed.

\section{The new idea: do not penalize---make the choice impossible
from the start}

The level-shift method ``penalizes the forbidden direction whenever it
is chosen.'' The rectangular-transformation approach changes the point
of view:

\begin{quote}
Eliminate the forbidden directions from the variables from the very
beginning.
\end{quote}

The variables the SCF varies are the MO coefficients \(\Cmat\). So,
instead of varying \(\Cmat\) directly and freely, we express it through
a smaller coefficient matrix \(\Cr\):
\[
  \Cmat = \Tmat\Cr .
\]

Here \(\Tmat\) is a tall \(N\times(N-m)\) matrix. Since the number of
columns is \(N-m\) rather than \(N\), the number of degrees of freedom
is reduced by \(m\). This is why we call it a ``rectangular
transformation.''

\(\Tmat\) is constructed so as to satisfy
\[
  \boxed{\Amat^\dagger\Smat\Tmat=\mathbf 0}
\]
Then
\begin{align*}
  \Amat^\dagger\Smat\Cmat
  &=
  \Amat^\dagger\Smat\Tmat\Cr \\
  &=
  \mathbf 0\,\Cr \\
  &=
  \mathbf 0 .
\end{align*}
This is crucially important. Whatever value \(\Cr\) takes, the
resulting \(\Cmat=\Tmat\Cr\) always satisfies
\[
  \Amat^\dagger\Smat\Cmat=\mathbf 0 .
\]
That is, no matter how the SCF proceeds, the unwanted hybrid orbitals
are unusable from the start.

\section{What does \(\Tmat\) represent?}

\(\Tmat\) is the matrix whose columns are the allowed directions only.
The condition characterizing the directions to be removed is
\[
  \Amat^\dagger\Smat\mathbf v=\mathbf 0 .
\]
A vector \(\mathbf v\) satisfying this is orthogonal, in the
\(\Smat\) metric, to the hybrid orbitals to be removed.

The space of all such \(\mathbf v\) is written
\[
  \ker(\Amat^\dagger\Smat) ,
\]
where \(\ker\) stands for the kernel, i.e., the null space.

The columns of \(\Tmat\) form a basis spanning this kernel. In other
words, \(\Tmat\) is ``the matrix whose columns are exactly the
directions orthogonal to the unwanted hybrid orbitals.''

If the dimension of the full AO space is \(N\) and \(m\) hybrid
orbitals are removed, the dimension of the remaining admissible space
is
\[
  N-m .
\]
Hence \(\Tmat\) is an \(N\times(N-m)\) matrix.

\section{An intuitive analogy}

The difference between the two methods is easy to grasp through the
following picture.

\begin{center}
\begin{tabular}{|p{0.22\linewidth}|p{0.35\linewidth}|p{0.33\linewidth}|}
\hline
\textbf{Method} & \textbf{Analogy} & \textbf{What happens} \\
\hline
Level shift & Build a steep hill on the forbidden road &
If the hill is finite, some intrusion remains possible \\
\hline
Rectangular transformation & Erase the road from the map &
That direction simply cannot be taken \\
\hline
\end{tabular}
\end{center}

The level shift is a ``penalty'' method. The rectangular transformation
is a ``remove it from the options'' method. Because of this difference,
the unwanted components are zero in principle under the rectangular
transformation.

\section{With the rectangular transformation, the level-shift term
vanishes}

The level-shifted matrix was
\[
  \widetilde{\Hcore}
  =
  \Hcore+\Venv+\Smat\Amat\Bmat\Amat^\dagger\Smat .
\]
Let us transform it into the admissible space, i.e., multiply by
\(\Tmat^\dagger\) from the left and \(\Tmat\) from the right:
\begin{align*}
\Tmat^\dagger\widetilde{\Hcore}\Tmat
&=
\Tmat^\dagger(\Hcore+\Venv)\Tmat \\
&\quad+
\Tmat^\dagger\Smat\Amat\Bmat\Amat^\dagger\Smat\Tmat .
\end{align*}
The second term can be rewritten as
\begin{align*}
&\Tmat^\dagger\Smat\Amat\Bmat\Amat^\dagger\Smat\Tmat \\
&\quad=
(\Amat^\dagger\Smat\Tmat)^\dagger
\Bmat
(\Amat^\dagger\Smat\Tmat).
\end{align*}
But \(\Tmat\) is constructed so that
\[
  \Amat^\dagger\Smat\Tmat=\mathbf 0 .
\]
Therefore
\[
  (\Amat^\dagger\Smat\Tmat)^\dagger
  \Bmat
  (\Amat^\dagger\Smat\Tmat)
  =
  \mathbf 0 .
\]

Hence the matrix used in the reduced space is
\[
  \boxed{
  \mathbf H_{\mathrm r}
  =
  \Tmat^\dagger(\Hcore+\Venv)\Tmat
  }
\]
and \(B\) no longer appears anywhere.

This is no accident. The role that \(B\) played---excluding the
unwanted directions---has been taken over entirely by the condition on
\(\Tmat\),
\[
  \Amat^\dagger\Smat\Tmat=\mathbf 0 .
\]

\section{The SCF equations in the reduced space}

The ordinary Roothaan--Hall equation is a generalized eigenvalue
problem because the AOs are non-orthogonal. After the rectangular
transformation, one solves the following equation in the reduced space:
\[
  \Fr\Cr=\Sr\Cr\boldsymbol\varepsilon ,
\]
where
\[
  \Sr=\Tmat^\dagger\Smat\Tmat
\]
and
\[
  \Fr
  =
  \Tmat^\dagger
  \left[
    \Hcore+\Venv+\Gmat[\Dmat]
  \right]
  \Tmat .
\]
\(\Gmat[\Dmat]\) denotes the two-electron part, i.e., the Coulomb and
exchange terms.

The density matrix is also mapped back from the reduced space to the
original AO space before use:
\[
  \Dmat = \Tmat\Dr\Tmat^\dagger .
\]
In this way, the two-electron integrals and the Fock-matrix
construction are handled as quantities in the original AO space, while
only the SCF degrees of freedom are reduced to dimension \(N-m\).

\section{\(S\)-orthonormalization}

It suffices for the columns of \(\Tmat\) to be orthogonal, in the
\(\Smat\) sense, to the orbitals to be removed. As constructed,
however, the overlap matrix of the reduced space,
\[
  \Sr=\Tmat^\dagger\Smat\Tmat ,
\]
need not be the identity.

One can therefore orthonormalize the columns of \(\Tmat\) in the
\(\Smat\) metric so that
\[
  \Tmat^\dagger\Smat\Tmat=\mathbf I .
\]
In that case the SCF equation in the reduced space becomes an ordinary
eigenvalue problem:
\[
  \Fr\Cr=\Cr\boldsymbol\varepsilon .
\]

In an implementation, one first finds a basis of
\[
  \ker(\Amat^\dagger\Smat)
\]
and then applies L\"owdin or canonical orthogonalization to it to
construct \(\Tmat\).

\section{A concrete example: an \(sp^3\) carbon}

Simplify the valence AOs of carbon to just the four functions
\[
  s,\quad p_x,\quad p_y,\quad p_z .
\]
From these four one can build four \(sp^3\) hybrid orbitals.

In the ideal case they can be written, for example, as
\begin{align*}
h_1 &= \frac12(s+p_x+p_y+p_z), \\
h_2 &= \frac12(s+p_x-p_y-p_z), \\
h_3 &= \frac12(s-p_x+p_y-p_z), \\
h_4 &= \frac12(s-p_x-p_y+p_z).
\end{align*}

If the severed bond points along \(h_1\), then \(h_1\) must not be used
in the calculation of fragment \(x\). The allowed directions are then
\(h_2,h_3,h_4\). Locally, a dimensional reduction
\[
  4\to 3
\]
takes place.

In this simple orthogonal-basis example, the transformation matrix on
the carbon is
\[
\Tmat_{\mathrm C}
=
\begin{array}{c|ccc}
 & h_2 & h_3 & h_4\\
\hline
s   & \frac12 & \frac12  & \frac12 \\
p_x & \frac12 & -\frac12 & -\frac12 \\
p_y & -\frac12 & \frac12 & -\frac12 \\
p_z & -\frac12 & -\frac12 & \frac12
\end{array}
\]
This \(\Tmat_{\mathrm C}\) can be viewed as the matrix obtained by
dropping the \(h_1\) column and keeping only the remaining three.

In actual calculations the basis set is much larger and the AOs are
non-orthogonal, so the coefficients are not such clean values of
\(1/2\). In practice one constructs \(\Tmat\) using \(\Smat\) so that
\[
  \Amat^\dagger\Smat\Tmat=\mathbf 0 .
\]
Still, the idea is the same: ``exclude the one cut-direction hybrid and
build the MOs from the remaining allowed directions only.''

\section{Handling of nuclei and electron counts when a bond is severed}

Both the rectangular transformation and the level-shift projection
operate only on the electronic variational space, i.e., the MO
coefficients \(\Cmat\). The nucleus of the \(\alpha\) carbon itself is
never divided by these operations. Since this point is easy to
misunderstand, we summarize it in its own section.

\subsection{Nuclei are not divided}

Even when a covalent bond is severed at the \(\alpha\) carbon, the
carbon nucleus belongs, as a point charge, entirely to one of the
fragments. Its nuclear charge is never halved, nor apportioned
according to the number of hybrid orbitals. Of the four hybrids, only
one (\(h_1\)) is removed; this merely reduces by one the number of
directions the electrons can use, and neither the position of the
nucleus nor the nuclear charge \(Z_{\mathrm C}\) changes.

\begin{verbatim}
   nucleus  Z_C = +6   <- point charge; belongs wholly to fragment x
      |                    (never divided)
   s, px, py, pz       <- all four AOs remain on the carbon
      |
   only the single direction h1 is removed
   (neither the nucleus nor the AOs themselves)
\end{verbatim}

If the \(\alpha\) carbon is assigned to fragment \(x\), the pieces are
distributed as follows.

\begin{itemize}
\item The nuclear charge \(+Z_{\mathrm C}\) enters the core Hamiltonian
      \(\Hcore\) of fragment \(x\) (as nucleus--electron attraction).
\item The nuclei on the partner-fragment side enter the environmental
      electrostatic potential \(\Venv\).
\item The AOs on the carbon (\(s,p_x,p_y,p_z\) and their
      generalizations) all remain in the basis. What the rectangular
      transformation removes is only ``one direction'' within the space
      these AOs span.
\end{itemize}

In short, neither nuclei nor AOs are divided or deleted; what is
transformed is only the set of directions available to the electrons.

\subsection{The number of AOs does not decrease; what decreases is the
number of variational degrees of freedom}

This point is particularly prone to confusion, so we emphasize it. The
rectangular transformation does not physically reduce the number of AO
basis functions \(N\). The basis functions remain \(N\) in number, and
the two-electron integrals are computed in the full \(N\)-dimensional
space. What is reduced is the number of degrees of freedom the SCF may
vary in searching for the MOs: from \(N\) to \(N-m\). In terms of the
equations, it is simply that \(\Cr\) in \(\Cmat=\Tmat\Cr\) is
\((N-m)\)-dimensional.

\subsection{Which fragment keeps the two electrons of the severed bond?}

When one covalent bond is severed, the assignment of the electron pair
(two electrons) that occupied that bond must be decided separately.
Removing a hybrid orbital only removes a ``direction''; it does not by
itself determine the electron count. In the standard FMO procedure, the
electron pair belonging to the severed bond is assigned to one of the
fragments, and the treatment of the boundary carbon (the so-called
bond-detached atom) is made consistent with that choice.

An implementation must therefore keep the following three items
mutually consistent:

\begin{enumerate}
\item To which fragment the nucleus (and core electrons) of the
      \(\alpha\) carbon are assigned.
\item To which fragment the electron pair of the severed bond is
      assigned.
\item Which hybrid orbitals are removed from the variational space.
\end{enumerate}

The rectangular transformation \(\Tmat\) (or the level shift
\(\Bmat\)) is responsible only for the third item. The first and
second are fixed separately as part of the definition of the fragment
partitioning.

\section{Similarities to and differences from the \(6\to5\)
transformation of \(d\) functions}

In quantum chemistry one sometimes transforms the six Cartesian \(d\)
functions into the five spherical-harmonic \(d\) functions. This is
also a \(6\to5\) dimensional reduction, so formally it resembles the
rectangular transformation. The meaning, however, is different.

\begin{center}
\begin{tabular}{|p{0.18\linewidth}|p{0.36\linewidth}|p{0.36\linewidth}|}
\hline
\textbf{Aspect} & \textbf{Cartesian \(d\) to spherical \(d\)} &
\textbf{FMO boundary rectangular transformation} \\
\hline
Dimensions & \(6\to5\) & \(N\to N-m\) \\
\hline
What is dropped & The redundant \(s\)-like component mixed into the
\(d\) functions & The hybrid orbital along the severed bond \\
\hline
How the matrix is determined & An essentially universal transformation
fixed by angular momentum & Depends on the bond direction, \(\Amat\),
and \(\Smat\) \\
\hline
Purpose & Tidy up the representation of the basis functions & Remove
unwanted components at the fragment boundary \\
\hline
\end{tabular}
\end{center}

Both reduce dimensions with a rectangular matrix, and in that respect
they are alike. However, the \(\Tmat\) of the FMO boundary treatment is
a matrix that imposes a physical constraint determined system by
system.

\section{Relation to the level-shift method}

The rectangular-transformation approach is not an unrelated alternative
to the level-shift method. Rather, it can be understood as follows:
\[
  \boxed{\text{deflation}\ \Longleftrightarrow\ B\to\infty}
\]
That is, the rectangular transformation corresponds to the
\(B\to\infty\) limit of the level-shift method.

In the level-shift method, the larger \(B\) is, the smaller the
components of the unwanted orbitals become. Ideally, taking \(B\) to
infinity would remove them completely. But an infinite \(B\) cannot be
used in actual numerical computation, and using a very large finite
value worsens the condition number.

The rectangular transformation realizes this ``ideal state of infinite
\(B\)'' without actually using \(B\). As a result, the tuning problem
of ``which value to choose for the level shift \(B\)'' disappears
altogether.

\section{Where do the matrices become non-square?}

A natural question at this point is: ``In the original approach with
the level shift \(B\), do the row and column dimensions also change,
making the matrices non-square?'' The short answer is:

\begin{quote}
In the level-shift method every matrix remains square,
\(N\times N\). Rectangular matrices appear only in the
rectangular-transformation approach, and even there the rectangular
one is the transformation matrix \(\Tmat\) itself; the matrix that is
finally diagonalized is again square.
\end{quote}

\subsection{Level-shift method: the shape stays \(N\times N\)}

The level-shifted matrix
\[
  \widetilde{\Hcore}
  =
  \Hcore+\Venv+\Smat\Amat\Bmat\Amat^\dagger\Smat
\]
is an \(N\times N\) square matrix. Tracking the sizes of the projection
term \(\Smat\Amat\Bmat\Amat^\dagger\Smat\),
\[
  \underbrace{\Smat}_{N\times N}\,
  \underbrace{\Amat}_{N\times m}\,
  \underbrace{\Bmat}_{m\times m}\,
  \underbrace{\Amat^\dagger}_{m\times N}\,
  \underbrace{\Smat}_{N\times N}
  \;\longrightarrow\; N\times N .
\]
Since \(m\) appears in the middle, the rank is \(m\), but the
``shape''---rows \(\times\) columns---remains \(N\times N\).

In the level-shift method, therefore:

\begin{itemize}
\item The generalized eigenvalue problem to be solved stays
      \(N\times N\); the dimension is not reduced.
\item \(N\) MOs are still obtained. The unwanted hybrid-orbital
      directions remain among the \(N\) solutions as high-lying virtual
      orbitals whose energies have been raised by \(B\). They do not
      disappear with a reduction of dimension; they are merely ``pushed
      up and left unoccupied.''
\end{itemize}

\subsection{Rectangular transformation: only \(\Tmat\) is rectangular}

In the rectangular-transformation approach, the transformation matrix
\[
  \Tmat:\quad N\times(N-m)\qquad(\text{rectangular})
\]
is the only rectangular matrix. Sandwiching with it, the transformed
matrices
\[
  \Fr=\Tmat^\dagger\Fmat\Tmat:\ (N-m)\times(N-m),\qquad
  \Sr=\Tmat^\dagger\Smat\Tmat:\ (N-m)\times(N-m)
\]
are again square, though their size has shrunk from \(N\) to
\(N-m\).

In other words, a ``rectangular matrix with different row and column
dimensions'' appears only as the \(\Tmat\) of the rectangular
transformation; the Fock and overlap matrices actually diagonalized are
square in both methods.

\begin{center}
\begin{tabular}{|p{0.26\linewidth}|p{0.3\linewidth}|p{0.3\linewidth}|}
\hline
\textbf{Matrix} & \textbf{Level-shift projection} &
\textbf{Deflation by rectangular transformation} \\
\hline
Transformation matrix & None (effectively the identity \(\mathbf I\),
\(N\times N\)) &
\(\Tmat\): \(N\times(N-m)\) (rectangular) \\
\hline
Matrix diagonalized & \(\widetilde{\Hcore}\): \(N\times N\) (square) &
\(\Fr,\Sr\): \((N-m)\times(N-m)\) (square) \\
\hline
Number of MOs & \(N\) & \(N-m\) \\
\hline
Treatment of unwanted directions & Remain among the \(N\) solutions at
high energy & Simply absent from the space \\
\hline
\end{tabular}
\end{center}

\subsection{Another viewpoint: the column count of \(\Tmat\) is the
answer}

That \(\Tmat\) is rectangular is itself the expression of the fact that
\(m\) directions have been dropped. A square \(N\times N\) matrix is a
transformation that ``keeps all \(N\) directions,'' and the level shift
corresponds to exactly this (the identity transformation plus a
projection term). A rectangular \(N\times(N-m)\) matrix is a
transformation that ``keeps only \(N-m\) directions,'' and this is the
essence of deflation.

The correspondence of the previous section,
\[
  \text{deflation}\ \Longleftrightarrow\ B\to\infty ,
\]
can be restated in these terms as well. In the \(B\to\infty\) limit,
\(m\) directions of the level shift's \(N\times N\) space are raised
infinitely high and become effectively unusable, and the remaining
\(N-m\) dimensions coincide exactly with the reduced space of the
rectangular transformation.

\section{Comparison of the two methods}

\begin{center}
\begin{tabular}{|p{0.18\linewidth}|p{0.36\linewidth}|p{0.36\linewidth}|}
\hline
\textbf{Item} & \textbf{Level-shift projection} &
\textbf{Deflation by rectangular transformation} \\
\hline
Basic idea & Penalize unwanted directions energetically & Remove
unwanted directions from the variables \\
\hline
Condition used & Add \(\Smat\Amat\Bmat\Amat^\dagger\Smat\) to the Fock
matrix & Use a \(\Tmat\) satisfying
\(\Amat^\dagger\Smat\Tmat=\mathbf 0\) \\
\hline
\(B\) & Required & Not required \\
\hline
Exactness of removal & Approximate for finite \(B\) & Exact in
principle \\
\hline
Residual components & Roughly \(O(B^{-1})\) & Zero to machine
precision \\
\hline
Dimension & Stays \(N\) & \(N-m\) \\
\hline
Numerical stability & Condition number can deteriorate for large
\(B\) & No \(B\)-induced conditioning problems \\
\hline
Viewpoint & Push out by making \(B\) large & Realize the
\(B\to\infty\) limit with a finite transformation \\
\hline
\end{tabular}
\end{center}

\section{Implementation outline}

An implementation proceeds as follows.
\begin{enumerate}
\item Build the cut-direction hybrid orbitals as AO coefficients:
      \(\Amat\).
\item Prepare the AO overlap matrix: \(\Smat\).
\item Find a basis spanning the kernel:
      \(\ker(\Amat^\dagger\Smat)\).
\item If desired, \(\Smat\)-orthonormalize that basis to obtain a
      \(\Tmat\) satisfying \(\Tmat^\dagger\Smat\Tmat=\mathbf I\).
\item Transform the Fock and overlap matrices into the reduced space:
      \[
        \Fr=\Tmat^\dagger\Fmat\Tmat,\qquad
        \Sr=\Tmat^\dagger\Smat\Tmat .
      \]
\item Solve the SCF equation in the reduced space:
      \[
        \Fr\Cr=\Sr\Cr\boldsymbol\varepsilon .
      \]
\item Map back to the original AO space as needed:
      \[
        \Cmat=\Tmat\Cr,\qquad
        \Dmat=\Tmat\Dr\Tmat^\dagger .
      \]
\end{enumerate}
Nowhere in this procedure does the level shift \(B\) appear.

\section{Summary}

In FMO, fragments are created by severing covalent bonds, so the hybrid
orbitals pointing along the severed bonds must be removed from the SCF
of the corresponding fragment.

The conventional level-shift method assigns a large energy \(B\) to the
unwanted hybrid orbitals to expel them from the occupied space. With
finite \(B\), however, the unwanted components never vanish completely,
and making \(B\) too large causes numerical instability.

In the rectangular-transformation approach, one constructs a matrix
\(\Tmat\) that spans only the space \(\Smat\)-orthogonal to the
unwanted hybrid orbitals, and sets
\[
  \Cmat=\Tmat\Cr .
\]
Since
\[
  \Amat^\dagger\Smat\Tmat=\mathbf 0 ,
\]
we have, for any \(\Cr\),
\[
  \Amat^\dagger\Smat\Cmat=\mathbf 0 .
\]
Therefore the unwanted hybrid orbitals never mix into the MOs in the
first place.

This method realizes the \(B\to\infty\) limit of the level-shift method
with finite matrix operations and without using \(B\). As a result, the
need to tune the level-shift value disappears, and the trade-off
between residual components and conditioning problems is avoided.

\section*{List of symbols}

\begin{longtable}{|p{0.28\linewidth}|p{0.62\linewidth}|}
\hline
\textbf{Symbol} & \textbf{Meaning} \\
\hline
\(\Hcore\) & Core Hamiltonian \\
\hline
\(\Venv\) & Environmental electrostatic potential from the other
fragments, etc. \\
\hline
\(Z_{\mathrm C}\) & Nuclear charge of the boundary (\(\alpha\)) carbon;
not divided, belongs to one fragment \\
\hline
\(\Smat\) & AO overlap matrix \\
\hline
\(\chi_\mu\) & AO basis function \\
\hline
\(\psi_p\) & MO \\
\hline
\(\Cmat\) & MO coefficients in the original AO space \\
\hline
\(\Cr\) & MO coefficients in the reduced space \\
\hline
\(\varphi_i^h\) & Hybrid orbital to be removed \\
\hline
\(\Amat\) & AO coefficients of the hybrid orbitals to be removed \\
\hline
\(m\) & Number of hybrid orbitals to be removed \\
\hline
\(N\) & Number of AO basis functions \\
\hline
\(\Bmat\) & Level-shift matrix \\
\hline
\(\Tmat\) & Rectangular transformation matrix spanning the admissible
space only \\
\hline
\(\Fmat\) & Fock matrix \\
\hline
\(\Gmat[\Dmat]\) & Two-electron term \\
\hline
\(\Dmat\) & Density matrix \\
\hline
\(\boldsymbol\varepsilon\) & Orbital energies \\
\hline
\end{longtable}

\end{document}
